% ----- CHAUPAI 21 -----

\newpage

\subsection*{\deva{एकविंशतितम चौपाई} \(Twenty-First Chaupai\)}
\addcontentsline{toc}{subsection}{\deva{एकविंशतितम चौपाई} \(Twenty-First Chaupai\) — \deva{राम दुआरे...}}

\begin{minipage}[t]{0.55\textwidth}
\vspace{0pt}

{\large \linespread{1.5}\selectfont
\deva{राम दुआरे तुम रखवारे।}\\
\deva{होत न आज्ञा बिनु पैसारे॥}
\par}

\vspace{0.5em}

{\large \linespread{1.5}\selectfont
\telu{రామ దుఆరే తుమ రఖవారే।}\\
\telu{హోత న ఆజ్ఞా బిను పైసారే॥}
\par}

\vspace{0.5em}

{\large \linespread{1.3}\selectfont
\textit{rāma duāre tuma rakhavāre |}\\
\textit{hota na ājñā binu paisāre ||}
\par}
\end{minipage}%
\hfill
\begin{minipage}[t]{0.42\textwidth}
\vspace{0pt}
\begin{tcolorbox}[anchorbox]
\textbf{The Gatekeeper of Discipline:} He guards the door to success (\textit{Ram Duaare Tum Rakhvare}). Nothing worthwhile in life is achieved without passing through the gates of discipline. Hanuman represents the hard work, routine, and discipline required before anyone can attain mastery or peace.
\vspace{0.5em}\newline Positioned at the very gates of the King, he became the ultimate symbol that success requires passing through strict discipline.\newline\null\hfill\href{https://www.valmiki.iitk.ac.in/sloka?field_kanda_tid=6&language=dv&field_sarga_value=128}{\textit{Yuddha Kanda, 128:88}}
\end{tcolorbox}
\end{minipage}

\vspace{0.5em}
\subsubsection*{\textbf{Visual Rhythm Map:}}
\begin{table}[h!]
\centering
\renewcommand{\arraystretch}{1.2}
\setlength{\tabcolsep}{0pt}
\begin{tabularx}{\textwidth}{|*{16}{>{\centering\arraybackslash\hsize=1.0\hsize}X|}}
\hline
\multicolumn{3}{|>{\centering\arraybackslash\hsize=3.0\hsize}X|}{\textbf{\guru{रा}\laghu{म}} \par (3)} &
\multicolumn{5}{>{\centering\arraybackslash\hsize=5.0\hsize}X|}{\textbf{\laghu{दु}\guru{आ}\guru{रे}} \par (5)} &
\multicolumn{2}{>{\centering\arraybackslash\hsize=2.0\hsize}X|}{\textbf{\laghu{तु}\laghu{म}} \par (2)} &
\multicolumn{6}{>{\centering\arraybackslash\hsize=6.0\hsize}X|}{\textbf{\laghu{र}\laghu{ख}\guru{वा}\guru{रे}} \par (6)} \\
\hline
\end{tabularx}
\vspace{-1.5em}
\end{table}

\begin{table}[h!]
\centering
\renewcommand{\arraystretch}{1.2}
\setlength{\tabcolsep}{0pt}
\begin{tabularx}{\textwidth}{|*{16}{>{\centering\arraybackslash\hsize=1.0\hsize}X|}}
\hline
\multicolumn{3}{|>{\centering\arraybackslash\hsize=3.0\hsize}X|}{\textbf{\guru{हो}\laghu{त}} \par (3)} &
\multicolumn{1}{>{\centering\arraybackslash\hsize=1.0\hsize}X|}{\textbf{\laghu{न}} \par (1)} &
\multicolumn{4}{>{\centering\arraybackslash\hsize=4.0\hsize}X|}{\textbf{\guru{आ}\guru{ज्ञा}} \par (4)} &
\multicolumn{2}{>{\centering\arraybackslash\hsize=2.0\hsize}X|}{\textbf{\laghu{बि}\laghu{नु}} \par (2)} &
\multicolumn{6}{>{\centering\arraybackslash\hsize=6.0\hsize}X|}{\textbf{\guru{पै}\guru{सा}\guru{रे}} \par (6)} \\
\hline
\end{tabularx}
\vspace{-1.5em}
\end{table}

\vspace{0.5em}
\textbf{ \deva{सरल-संस्कृतम्}:}\\
 \deva{श्रीरामस्य द्वारे भवान् एव रक्षकः (द्वारपालः) अस्ति।}\\
 \deva{भवतः आज्ञां विना तत्र प्रवेशः न भवति॥}\\

At the door of Lord Rama, you are the protector/guard.\\
Without your permission, entry is not possible.


\vspace{1em}
\newpage
\textbf{\deva{प्रतिपदार्थः} (Meaning of each word):}
\begin{table}[h!]
\centering
\renewcommand{\arraystretch}{1.2}
\begin{tabularx}{\textwidth}{@{} l l X X @{}}
\toprule
\textbf{अवधी} & \textbf{संस्कृतम्} & \textbf{English} & \textbf{Grammar \& Notes} \\
\midrule
\deva{राम दुआरे} / \telu{రామ దుఆరే}  &  \deva{रामस्य द्वारे}  &  At Rama's door  & `duaare` is locative \\
\deva{तुम} \gotchaicon / \telu{తుమ}  &  \deva{त्वम्}  &  You  &   \\
\deva{रखवारे} / \telu{రఖవారే}  &  \deva{रक्षकः}  &  Protector / Guard  & \\ 
\deva{होत न} / \telu{హోత న}  &  \deva{न भवति}  &  Does not happen  &   \\
\deva{आज्ञा} / \telu{ఆజ్ఞా}  &  \deva{आज्ञाम्}  &  Command / Permission  &   \\
\deva{बिनु} / \telu{బిను}  &  \deva{विना}  &  Without  &   \\
\deva{पैसारे} / \telu{పైసారే}  &  \deva{प्रवेशः}  &  Entry / Passage  & Awadhi noun \\
\bottomrule
\end{tabularx}
\end{table}

\begin{tcolorbox}[poetrybox]
\textbf{\deva{त्वम्} $\rightarrow$ \deva{तुम} \gotchaicon\ — The Eternal ''You'':}\\
Sanskrit \deva{त्वम्} (nominative ''you'') traces back to Proto-Indo-European \textit{*túh₂} — the same root that gives English ''thou,'' German \textit{du}, and Persian \textit{tu}. Awadhi condenses this to the unflinching \deva{तुम}, stripping away all case distinctions. Here the word carries exceptional weight: Tulsidas addresses Hanuman directly and personally, making this verse a powerful statement that \textit{only you} — not any other deity — stand guard at Rama's door. It is an intimate, almost devotional ''you.''
\end{tcolorbox}

\newpage
